\subsection{Justificativa da escolha}
O controlador PD foi adotado por conter uma estrutura simples e possui baixo custo computacional, além da ampla aplicação em sistemas de controle de atitude e em manipuladores robóticos, onde é empregado para estabilização e rastreamento de trajetória.
Embora técnicas mais avançadas, como o LQR e o controle robusto, possam oferecer melhor desempenho diante de incertezas e de perturbações, o controlador PD possui resposta adequada às condições consideradas neste trabalho, sendo possível avaliar o comportamento dinâmico do sistema de forma clara e controlada.